% !TeX program = xelatex
% ==========================================================================
% Звіт з лабораторної роботи — ДВВ «Основи прикладної фізики» (1.26-27), УКУ.
% Викладач: Володимир Коломієць.
% Це ЄДИНИЙ файл, який вам потрібно редагувати. Уся верстка й брендинг
% живуть у ucaphyslab.cls — не змінюйте шрифти, поля, кеглі й кольори.
%
% Компіляція: XeLaTeX + BibTeX (Overleaf робить це автоматично).
% Локально:   latexmk (див. README.md у корені репозиторію).
% ==========================================================================
\documentclass[ukrainian]{ucaphyslab}
% Опції класу (можна комбінувати через кому):
%   ukrainian | english  — мова документа й логотипу (за замовчуванням ukrainian)
%   slogan                — додає гасло УКУ в колонтитул титулки (за замовчуванням вимкнено)

% --------------------------------------------------------------------------
% Метадані титульної сторінки
% --------------------------------------------------------------------------
\labtitle{Математичний і фізичний (оборотний) маятники}
\labsubtitle{Лабораторна робота № 2}

\student{Довбенчук Олена}
\student{Кульбаба Андрій}
\student{Магарита Ія}
\student{Пастарнак Марко}

\teacher{Володимир Коломієць} % якщо рядок не потрібен — видаліть текст і залиште \teacher{}

% Підписи на титулці (замість «Виконав(ла)» / «Перевірив(ла)» з класу)
\makeatletter
\renewcommand{\ucu@labelstudent}{Студенти}
\renewcommand{\ucu@labelteacher}{Викладач}
\makeatother

\labdate{2026-10-13}

\usepackage{float}
\usepackage{pgfplots}
\pgfplotsset{compat=1.18}
\sisetup{uncertainty-mode=separate}
\definecolor{series1}{HTML}{2A78D6}
\definecolor{series2}{HTML}{EB6834}

\begin{document}
\maketitlepage

\section{Мета роботи}
\begin{enumerate}
  \item Визначити прискорення вільного падіння за періодом малих коливань математичного маятника та за зведеною довжиною оборотного фізичного маятника.
  \item Вивчити вільні коливання маятника, що задовольняє моделі математичного маятника, і вільні коливання оборотного маятника.
\end{enumerate}

\section{Прилади та обладнання}
\begin{itemize}
  \item \textbf{Установка} для вивчення коливань маятників: основа з регульованими опорами, колона зі шкалою, верхній кронштейн.
  \item \textbf{Математичний маятник}: металева кулька на біфілярному підвісі (дві нитки). Довжину підвісу змінюють гвинтом і зчитують зі шкали на колоні; у роботі \(l > \SI{50}{\centi\meter}\).
  \item \textbf{Фізичний (оборотний) маятник}: сталевий стержень, дві опорні призми (ножі), повернені лезами одна до одної, і два масивні диски. Ножі й диски переміщують уздовж стержня і фіксують. Відстань між ножами (зведена довжина) визначають за кільцевими заглибленнями на стержні з кроком \SI{10}{\milli\meter}.
  \item \textbf{Секундоміри} (час вимірювали 3 спостерігачі одночасно).
  \item \textbf{Рулетка} для вимірювання відстані між ножами.
  \item \textbf{Ваги} для вимірювання маси важка.
\end{itemize}


\section{Теоретичні відомості}
% Основні формули, з яких обчислюються шукані величини, з короткими
% поясненнями позначень.
Математичний маятник — матеріальна точка на невагомому нерозтяжному підвісі довжини \(l\). За малих коливань період і прискорення вільного падіння
\begin{equation}
  T = 2\pi\sqrt{\frac{l}{g}}, \qquad
  g = \frac{4\pi^{2} l}{T^{2}}.
  \tag{2.13}
  \label{eq:g-math}
\end{equation}

Оборотний маятник — фізичний маятник із двома паралельними осями підвісу. Якщо періоди коливань відносно цих осей однакові, відстань між ними дорівнює зведеній довжині \(l_{\text{зв}}\), і
\begin{equation}
  g = \frac{4\pi^{2} l_{\text{зв}}}{T^{2}}.
  \tag{2.33б}
  \label{eq:g-phys}
\end{equation}
\section{Хід роботи та результати вимірювань}
Час десяти повних коливань одночасно вимірювали три спостерігачі ручними секундомірами.
У таблицях \(t_1\), \(t_2\), \(t_3\) — їхні покази, \(t_\text{сер}\) — середнє, період \(T = t_\text{сер}/10\).

\begin{table}[H]
  \centering
  \begin{tabular}{@{}llc@{}}
    \toprule
    Величина & Позначення & Значення \\
    \midrule
    Довжина підвісу математичного маятника & \(l\) & \SI{66}{\centi\meter} \\
    Маса кульки & \(m\) & \SI{177}{\gram} \\
    Відстань між ножами & \(l_{\text{зв}}\) & \SI{64.3}{\centi\meter} \\
    Табличне значення & \(g\) & \SI{9.81}{\meter\per\second\squared} \\
    \bottomrule
  \end{tabular}
  \caption{Константи та параметри установки}
  \label{tab:constants}
\end{table}

\subsection{Математичний маятник}
\begin{enumerate}
  \item Встановили однакові довжини ниток біфілярного підвісу і виміряли довжину \(l = \SI{66}{\centi\meter}\) за шкалою на колоні.
  \item Відхилили кульку на \(4\)--\(5^{\circ}\) від рівноваги. Три спостерігачі одночасно виміряли секундомірами час десяти повних коливань.
  \item Вимір повторили 10 разів за тієї самої довжини підвісу.
\end{enumerate}

\begin{table}[H]
  \centering
  \begin{tabular}{@{}c ccc cc@{}}
    \toprule
    № & \(t_1\), с & \(t_2\), с & \(t_3\), с & \(t_\text{сер}\), с & \(T\), с \\
    \midrule
    1  & 16.38 & 16.81 & 16.42 & 16.537 & 1.654 \\
    2  & 16.41 & 16.79 & 16.38 & 16.527 & 1.653 \\
    3  & 16.83 & 16.64 & 16.45 & 16.640 & 1.664 \\
    4  & 16.65 & 16.81 & 16.65 & 16.703 & 1.670 \\
    5  & 16.48 & 16.73 & 16.37 & 16.527 & 1.653 \\
    6  & 16.58 & 16.68 & 16.37 & 16.543 & 1.654 \\
    7  & 16.40 & 16.39 & 16.29 & 16.360 & 1.636 \\
    8  & 16.46 & 16.79 & 16.37 & 16.540 & 1.654 \\
    9  & 16.60 & 16.86 & 16.32 & 16.593 & 1.659 \\
    10 & 16.39 & 16.77 & 16.35 & 16.503 & 1.650 \\
    \bottomrule
  \end{tabular}
  \caption{Математичний маятник: час десяти коливань і період (\(l = \SI{66}{\centi\meter}\)). Середній період \(T_\text{сер} = \SI{1.655}{\second}\)}
  \label{tab:math}
\end{table}

\subsection{Фізичний (оборотний) маятник}
\begin{enumerate}
  \item Диски зафіксували на стержні несиметрично: один поблизу кінця, другий поблизу середини. Ножі закріпили лезами до середини стержня: ніж №1 поблизу вільного кінця, ніж №2 між дисками.
  \item Маятник по черзі підвішували на кожен ніж, відхиляли на \(4\)--\(5^{\circ}\) і трьома секундомірами вимірювали час десяти коливань.
  \item Ніж №2 переміщували вздовж стержня, доки періоди на двох ножах не зблизились. Відстань між ножами виміряли рулеткою: \(l_{\text{зв}} = \SI{64.3}{\centi\meter}\). У цьому положенні виконали 5 пар вимірів.
\end{enumerate}

\begin{table}[H]
  \centering
  \small
  \setlength{\tabcolsep}{4pt}
  \begin{tabular}{@{}c ccc c ccc c c@{}}
    \toprule
    & \multicolumn{4}{c}{Ніж №1} & \multicolumn{4}{c}{Ніж №2} & \\
    \cmidrule(lr){2-5} \cmidrule(lr){6-9}
    № & \(t_1\), с & \(t_2\), с & \(t_3\), с & \(T_1\), с
      & \(t_1\), с & \(t_2\), с & \(t_3\), с & \(T_2\), с
      & \(|T_1-T_2|\), с \\
    \midrule
    1 & 15.88 & 16.09 & 15.80 & 1.592 & 15.94 & 16.22 & 16.02 & 1.606 & 0.014 \\
    2 & 16.07 & 16.11 & 16.15 & 1.611 & 15.96 & 16.30 & 16.12 & 1.613 & 0.002 \\
    3 & 16.07 & 16.20 & 15.74 & 1.600 & 15.94 & 16.24 & 15.85 & 1.601 & 0.001 \\
    4 & 15.80 & 16.21 & 16.02 & 1.601 & 15.85 & 16.12 & 16.00 & 1.599 & 0.002 \\
    5 & 15.94 & 16.40 & 16.05 & 1.613 & 15.89 & 16.00 & 15.89 & 1.593 & 0.020 \\
    \bottomrule
  \end{tabular}
  \caption{Оборотний маятник: час десяти коливань на двох ножах (\(l_{\text{зв}} = \SI{64.3}{\centi\meter}\))}
  \label{tab:phys}
\end{table}

\section{Обробка результатів}
Прискорення вільного падіння обчислено за середніми періодами: для математичного маятника за \eqref{eq:g-math}, для оборотного — за \eqref{eq:g-phys}.

\subsection{Математичний маятник}
За десятьма періодами табл.~\ref{tab:math} \(T_\text{сер} = \SI{1.655}{\second}\), \(l = \SI{0.66}{\meter}\):
\begin{equation}
  g = \frac{4\pi^{2}\cdot \num{0.66}}{(\num{1.655})^{2}} = \SI{9.51}{\meter\per\second\squared}.
  \label{eq:g-math-num}
\end{equation}

\subsection{Оборотний маятник}
Середній період на ножі №1 дорівнює \SI{1.603}{\second}, на ножі №2 — \SI{1.602}{\second}. Різниця \SI{0.001}{\second} менша за \(0{,}5\,\%\), тож для середніх умова \(T_1 \approx T_2\) виконана. Середнє за всіма десятьма вимірами \(T_\text{сер} = \SI{1.603}{\second}\), \(l_{\text{зв}} = \SI{0.643}{\meter}\):
\begin{equation}
  g = \frac{4\pi^{2}\cdot \num{0.643}}{(\num{1.603})^{2}} = \SI{9.88}{\meter\per\second\squared}.
  \label{eq:g-phys-num}
\end{equation}

\subsection{Похибки}
Випадкова похибка періоду — довірча похибка середнього (\(P = 0{,}95\), \(t_{0.95;\,9} = 2.26\), \(n = 10\)):
\begin{equation*}
  \Delta T_\text{вип} = t_{0.95;\,9}\sqrt{\frac{\sum(T_i - T_\text{сер})^{2}}{n(n-1)}}.
\end{equation*}
Систематична похибка періоду — ціна поділки секундоміра \SI{0.01}{\second}, поділена на десять коливань. Звідси \(\Delta T_\text{сист} = \SI{0.001}{\second}\). Вона не зменшується з числом спроб.

Довжину математичного маятника записано до \SI{1}{\centi\meter}, тому \(\Delta l = \SI{0.5}{\centi\meter}\). Відстань між ножами виміряно рулеткою з ціною поділки \SI{1}{\milli\meter}, тому \(\Delta l_{\text{зв}} = \SI{0.5}{\milli\meter}\).

Для \(g = 4\pi^{2} l/T^{2}\) випадкова похибка походить від розкиду періоду, систематична — від приладів:
\begin{align*}
  \Delta g_\text{вип} &= g\cdot\frac{2\Delta T_\text{вип}}{T}, \\
  \Delta g_\text{сист} &= g\sqrt{\left(\frac{\Delta l}{l}\right)^{2} + \left(\frac{2\Delta T_\text{сист}}{T}\right)^{2}}, \\
  \Delta g &= \sqrt{\Delta g_\text{вип}^{2} + \Delta g_\text{сист}^{2}}.
\end{align*}

\begin{table}[H]
  \centering
  \begin{tabular}{@{}lcc@{}}
    \toprule
    Величина & Математичний & Оборотний \\
    \midrule
    \(T_\text{сер}\), с & \num{1.655} & \num{1.603} \\
    \(\Delta T_\text{вип}\), с & \num{0.0064} & \num{0.0055} \\
    \(\Delta T_\text{сист}\), с & \num{0.001} & \num{0.001} \\
    \midrule
    \(g\), м/с\(^2\) & \num{9.51 \pm 0.10} & \num{9.88 \pm 0.07} \\
    \(\Delta g_\text{вип}\), м/с\(^2\) & \num{0.07} & \num{0.07} \\
    \(\Delta g_\text{сист}\), м/с\(^2\) & \num{0.07} & \num{0.01} \\
    Довірчий інтервал, м/с\(^2\) & \([9.41;\ 9.61]\) & \([9.81;\ 9.95]\) \\
    \midrule
    Табличне \(g\), м/с\(^2\) & \num{9.81} & \num{9.81} \\
    Відносна похибка \(\Delta g/g\), \% & \num{1.1} & \num{0.7} \\
    Відхилення від табличного, \% & \num{-3.1} & \num{0.7} \\
    \bottomrule
  \end{tabular}
  \caption{Прискорення вільного падіння та похибки (\(P = 0.95\))}
  \label{tab:g}
\end{table}

\begin{figure}[H]
  \centering
  \begin{tikzpicture}
    \begin{axis}[
      width=0.72\linewidth, height=6.2cm,
      ybar, bar width=28pt,
      xlabel={}, ylabel={$g$, м/с$^2$},
      xmin=0.4, xmax=2.6, xtick={1,2},
      xticklabels={Математичний, Оборотний},
      ymin=9.2, ymax=10.15,
      ymajorgrids, grid style={gray!25},
      axis lines*=left,
      error bars/y dir=both, error bars/y explicit,
      legend style={draw=none, at={(0.5,1.02)}, anchor=south, font=\small},
    ]
      \addplot[fill=series1, draw=none] coordinates {
        (1,9.51) +- (0,0.10)
        (2,9.88) +- (0,0.07)
      };
      \addlegendentry{експеримент}
      \addplot[color=series2, dashed, thick, forget plot] coordinates {(0.4,9.81) (2.6,9.81)};
    \end{axis}
  \end{tikzpicture}
  \caption{Виміряне \(g\) з повною похибкою. Штрихова лінія — табличне значення \SI{9.81}{\meter\per\second\squared}}
  \label{fig:g}
\end{figure}

Табличне значення для математичного маятника лежить поза довірчим інтервалом: виміряний період \SI{1.655}{\second} більший за \(2\pi\sqrt{l/g} = \SI{1.63}{\second}\), і цей зсув перевищує розкид між спробами. Для оборотного маятника табличне значення лежить на нижній межі інтервалу, відхилення \(0{,}7\,\%\) не перевищує похибки вимірювання.

Можливі причини зсуву для математичного маятника: час відлічували вручну, тож спільна затримка реакції збільшує період і не входить у розкид між спробами; довжину зчитано до сантиметра, і вона може не збігатися з відстанню від точки підвісу до центра кульки.

\section{Висновки}
Прискорення вільного падіння визначено за періодом малих коливань математичного маятника та за зведеною довжиною оборотного маятника. Результати порівняно з табличним значенням \(g = \SI{9.81}{\meter\per\second\squared}\).

Для математичного маятника \(g = (9{,}51 \pm 0{,}10)\)~м/с\(^2\). Табличне значення лежить поза довірчим інтервалом \([9{,}41;\ 9{,}61]\): відхилення \(-3{,}1\,\%\) більше за відносну похибку вимірювання \(1{,}1\,\%\). Для оборотного маятника \(g = (9{,}88 \pm 0{,}07)\)~м/с\(^2\), інтервал \([9{,}81;\ 9{,}95]\). Табличне значення лежить на нижній межі інтервалу, відхилення \(0{,}7\,\%\) не перевищує похибки вимірювання.

Зсув для математичного маятника дає завищений період. Час відлічували вручну, тож спільна затримка реакції збільшує \(T\) і не входить у розкид між спробами. Довжину зчитано до сантиметра, і вона може не збігатися з відстанню від точки підвісу до центра кульки.

% Розкоментуйте, якщо цитуєте джерела через \cite{} (записи — у references.bib):
% \ucuprintbibliography

\end{document}
